% Part: first-order-logic
% Chapter: completeness
% Section: completeness-theorem

\documentclass[../../../include/open-logic-section]{subfiles}

\begin{document}

\iftag{FOL}
      {\olfileid{fol}{com}{cth}}
      {\olfileid{pl}{com}{cth}}

\olsection{પૂર્ણતા પ્રમેય}

\begin{explain}
આપણાં પરિણામો હવે ભેગાં કરીએ: તેમાંથી પૂર્ણતા પ્રમેય મળે છે.
\end{explain}

\begin{thm}[પૂર્ણતા પ્રમેય]
\ollabel{thm:completeness}
ધારો કે~$\Gamma$ !!{sentence}sનો ગણ છે. જો~$\Gamma$ સુસંગત હોય, તો તે
સંતોષ્ય છે.
\end{thm}

\begin{proof}
ધારો કે~$\Gamma$ સુસંગત છે.\iftag{FOL}{ \olref[hen]{lem:henkin} મુજબ
સંતૃપ્ત સુસંગત ગણ $\Gamma'\supseteq\Gamma$ છે.}{}
\olref[lin]{lem:lindenbaum} મુજબ એવો
$\Gamma^*\supseteq\iftag{FOL}{\Gamma'}{\Gamma}$ છે જે સુસંગત અને
!!{complete} છે.
\iftag{FOL}{%
  $\Gamma'\subseteq\Gamma^*$ હોવાથી દરેક !!{formula}~$!A(x)$ માટે
  ~ $\Gamma^*$માં આ સ્વરૂપનું !!a{sentence} હોય છે:
  \iftag{prvEx}
        {$\lexists[x][!A(x)] \lif !A(c)$}
        {$\lnot\lforall[x][!A(x)] \lif \lnot !A(c)$}.
  તેથી~$\Gamma^*$ સંતૃપ્ત છે. જો~$\Gamma$માં~$\eq$ ન આવતું હોય, તો}{તો}
\olref[mod]{lem:truth} મુજબ
$\iftag{FOL}{\Sat{M(\Gamma^*)}}{\pSat{v(\Gamma^*)}}{!A}$ જો અને માત્ર
જો $!A\in\Gamma^*$. ખાસ કરીને તેમાંથી દરેક $!A\in\Gamma$ માટે
$\iftag{FOL}{\Sat{M(\Gamma^*)}}{\pSat{v(\Gamma^*)}}{!A}$ ફલિત થાય છે,
તેથી~$\Gamma$ સંતોષ્ય છે.\iftag{FOL}{%
  જો~$\Gamma$માં~$\eq$ આવતું હોય, તો \olref[ide]{lem:truth} મુજબ દરેક
  !!{sentence}sમાંના~$!A$ માટે $\Sat{\equivclass{M}{\approx}}{!A}$ જો અને
  માત્ર જો $!A\in\Gamma^*$. ખાસ કરીને દરેક $!A\in\Gamma$ માટે
  $\Sat{\equivclass{M}{\approx}}{!A}$, તેથી~$\Gamma$ સંતોષ્ય છે.}{}
\end{proof}

\begin{cor}[પૂર્ણતા પ્રમેય, બીજું નિરૂપણ]
\ollabel{cor:completeness}
દરેક~$\Gamma$ અને !!{sentence}sમાંના~$!A$ માટે: જો $\Gamma\Entails !A$, તો
$\Gamma\Proves !A$.
\end{cor}

\begin{proof}
નોંધો કે \olref{cor:completeness} અને \olref{thm:completeness}માં
~$\Gamma$ પર સર્વવાચક પરિમાણ છે. ગૂંચવણ ન થાય તે માટે
\olref{thm:completeness}ને જુદા ચલથી ફરી કહીએ: !!{sentence}sના કોઈ પણ
ગણ~$\Delta$ માટે, જો~$\Delta$ સુસંગત હોય, તો તે સંતોષ્ય છે.
પ્રતિપક્ષન મુજબ, જો~$\Delta$ સંતોષ્ય ન હોય, તો~$\Delta$ અસુસંગત છે. આનો
ઉપપ્રમેય સાબિત કરવા ઉપયોગ કરીશું.

ધારો કે $\Gamma\Entails !A$. ત્યારે \olref[syn][sem]{prop:entails-unsat}
મુજબ $\Gamma\cup\{\lnot !A\}$ અસંતોષ્ય છે. આપણા~$\Delta$ તરીકે
$\Gamma\cup\{\lnot !A\}$ લઈએ, તો \olref{thm:completeness}નું પહેલું
નિરૂપણ કહે છે કે $\Gamma\cup\{\lnot !A\}$ અસુસંગત છે. હવે
\iftag{FOL}{%
  \tagrefs{prfAX/{fol:axd:prv:prop:prov-incons},
    prfSC/{fol:seq:prv:prop:prov-incons},
    prfND/{fol:ntd:prv:prop:prov-incons},
    prfTab/{fol:tab:prv:prop:prov-incons}}}{%
  \tagrefs{prfAX/{pl:axd:prv:prop:prov-incons},
    prfSC/{pl:seq:prv:prop:prov-incons},
    prfND/{pl:ntd:prv:prop:prov-incons},
    prfTab/{pl:tab:prv:prop:prov-incons}}}
મુજબ $\Gamma\Proves !A$.
\end{proof}

\tagprob{FOL}
\begin{prob}
\olref[fol][com][cth]{cor:completeness} વાપરીને
\olref[fol][com][cth]{thm:completeness} સાબિત કરો અને આ રીતે પૂર્ણતા
પ્રમેયનાં બંને નિરૂપણ સમતુલ્ય છે તેમ બતાવો.
\end{prob}
\tagendprob

\tagprob{notFOL}
\begin{prob}
\olref[pl][com][cth]{cor:completeness} વાપરીને
\olref[pl][com][cth]{thm:completeness} સાબિત કરો અને આ રીતે પૂર્ણતા
પ્રમેયનાં બંને નિરૂપણ સમતુલ્ય છે તેમ બતાવો.
\end{prob}
\tagendprob

\tagprob{FOL}
\begin{prob}
કોઈ !!a{derivation} તંત્ર પૂર્ણ હોય તે માટે તેના નિયમો એટલા સમર્થ હોવા
જોઈએ કે તે દરેક અસંતોષ્ય ગણને અસુસંગત સાબિત કરી શકે. પૂર્ણતા સાબિત કરવા
!!{derivation}ના કયા નિયમો જરૂરી હતા? આમાંથી કોઈ નિયમ સાબિતીમાં ક્યાંય
વપરાયો નથી? આ પ્રશ્નોના જવાબ માટે એવી યાદી અથવા આકૃતિ બનાવો જે બતાવે
કે \olref[fol][com][cth]{thm:completeness}ની સાબિતી સુધી લઈ જતાં કયાં
પરિણામોમાં !!{derivation}ના કયા નિયમો વપરાયા હતા. આ સાબિતીઓમાં નિયમોના
કોઈ અવ્યક્ત ઉપયોગ પણ નોંધવાનું ભૂલશો નહિ.
\end{prob}
\tagendprob

\tagprob{notFOL}
\begin{prob}
કોઈ !!a{derivation} તંત્ર પૂર્ણ હોય તે માટે તેના નિયમો એટલા સમર્થ હોવા
જોઈએ કે તે દરેક અસંતોષ્ય ગણને અસુસંગત સાબિત કરી શકે. પૂર્ણતા સાબિત કરવા
!!{derivation}ના કયા નિયમો જરૂરી હતા? આમાંથી કોઈ નિયમ સાબિતીમાં ક્યાંય
વપરાયો નથી? આ પ્રશ્નોના જવાબ માટે એવી યાદી અથવા આકૃતિ બનાવો જે બતાવે
કે \olref[pl][com][cth]{thm:completeness}ની સાબિતી સુધી લઈ જતાં કયાં
પરિણામોમાં !!{derivation}ના કયા નિયમો વપરાયા હતા. આ સાબિતીઓમાં નિયમોના
કોઈ અવ્યક્ત ઉપયોગ પણ નોંધવાનું ભૂલશો નહિ.
\end{prob}
\tagendprob

\end{document}
